\section{IMP-UI-004 — Notificações, deep link de roteiro e estado offline}

\textbf{Escopo.} Esta verificação cobre a caixa interna e a página dedicada de
notificações (UI-12), o comportamento de cache em memória/offline (UI-13) e o
deep link de roteiro exigido por RN-M13-04. Não há persistência de conteúdo de
notificação em \texttt{localStorage}; mutações continuam sendo confirmadas por
callable e validação server-side.

\textbf{Implementação verificada.} A página e o dropdown limpam a lista quando
o UID muda, encerram o listener anterior e preservam a última leitura autorizada
quando o listener perde a conexão. O estado usa \texttt{snapshot.metadata.fromCache}
e eventos \texttt{online/offline}; quando aplicável, a UI exibe o aviso
``Você está offline'' ou ``Notificações possivelmente desatualizadas''. O
entidade-alvo \texttt{Roteiro} aponta para
\texttt{/turmas?roteiros=1}, que abre o gerenciamento de roteiros e volta a
revalidar o acesso na tela protegida.

\textbf{Evidência E2E (Node.js 20).}

\begin{itemize}
  \item \texttt{12-notificacoes.spec.ts}: E2E-NOTIF-001--003, 3/3 verdes.
  \item \texttt{15-notificacoes-page.spec.ts}: UI-004-E2E-001 (deep link de
  roteiro), UI-004-E2E-002 (aviso offline) e regressão legada, 3/3 verdes.
\end{itemize}

As suítes foram executadas separadamente, cada uma com uma única preparação da
seed, pois compartilham o banco do emulador. A execução combinada dos dois
arquivos não é evidência válida: o primeiro arquivo deixa notificações de
arquivamento para a sequência do segundo.
